\begingroup\small
\begin{longtable}{l l p{8.0cm}}
\hline
part & type & statements \\
\hline
\endhead
Part I (triples) & proved & \textbf{I.2.1~height inequality}\newline \textbf{I.2.4~two wings}\newline \textbf{I.4.6~reduction to a witness statement}\newline \textbf{I.4.8~the divisor count of the Wieferich conditions}\newline \textbf{I.6.10~Wieferich primes of the unit and the primitive parts}\newline \emph{also:} I.2.2~the constant in the height bound; I.2.3~kernel bound; I.3.2~weaker bounds with less work; I.4.3~valuation identity; I.4.4~exponent-3 condition; I.4.7~the $T_1$-part of a lattice point; I.4.10~what the witness condition does not follow from; I.4.11~where a witness in an outer number can sit; I.4.12~two elementary constraints, and the size of the gap; I.4.13~the cyclotomic axis \\
 & computed & \textbf{I.3.1~reduction}\newline \textbf{I.4.5~the $b$-box}\newline \textbf{I.5.1~the strip}\newline \textbf{I.5.3~smooth kernels}\newline \textbf{I.5.4~square-type middles}\newline \textbf{I.6.2~no triple below $\Z{ohnetripelsicher}$, and the pairs at distance two with middle divisible by four}\newline \textbf{I.6.8~the primitive part over an explicit finite set}\newline \textbf{I.6.13~a second primitive prime factor}\newline \textbf{I.7.1~profile}\newline \emph{also:} I.5.5~the share of the candidates covered; I.6.3~ladders for the known kernels dividing $U_1$; I.6.4~the smallest witnesses; I.6.5~the condition modulo $36$ and the witnesses; I.6.6~how many conditions there are, and how often they are met; I.6.7~middles that are powers of two; I.6.11~the Wieferich events of the units; I.6.12~Lucas--Wieferich primes; I.8.2~Yokoi's representation; I.8.3~small units \\
 & measured & \textbf{I.6.1~exponents at the rank} \\
 & open & \emph{finiteness of the kernels of wing (ii)} (Problem~\problemvon{I:q:walsh}); \emph{the primitive part of $T_d(m)$} (Problem~\problemvon{I:q:primindex}); \emph{a middle that is a power of $2$} (Problem~\problemvon{I:q:zweierpotenz}); \emph{Reinhart's condition (C)} (Problem~\problemvon{I:q:reinhart}) \\
\hline
Part II (pairs) & proved & \textbf{II.6.2~the count is a sum over towers}\newline \textbf{II.6.4~localization of the pair count}\newline \textbf{II.6.6~where Tao's argument is weakest}\newline \emph{also:} II.2.1~powerful numbers; II.6.5~what the localization says, and what it does not say; II.7.1~pairs on one Pell equation \\
 & computed & \textbf{II.3.1~the known pairs}\newline \textbf{II.4.1~explicit lower bound}\newline \textbf{II.5.2~the families with a square overtake the comparison function}\newline \textbf{II.6.1~the middles that are not divisible by four}\newline \textbf{II.6.3~lower bound at every height}\newline \textbf{II.7.2~the third number}\newline \emph{also:} II.3.2~completeness inside the families \\
 & heuristic & \textbf{II.5.1~comparison function} \\
 & open & \emph{convergence of the rates} (Problem~\problemvon{II:q:konvergenz}); \emph{growth of the tower rate} (Problem~\problemvon{II:q:wachstum}); \emph{the kernels that divide $U_1$} (Problem~\problemvon{II:q:ezwei}) \\
\hline
Part III (further observations) & proved & \textbf{III.2.2~two explicit chains of powerful values of $\Kreis{3}$ and $\Kreis{6}$}\newline \textbf{III.2.5~the simplest case of $\Kreis{8}$} \\
 & computed & \textbf{III.1.2~the head of the gap spectrum}\newline \textbf{III.2.3~powerful values up to a bound}\newline \emph{also:} III.1.3~a ceiling for common divisors; III.2.1~the values $x^2 + 1$; III.2.4~a class without squares \\
 & measured & \textbf{III.1.4~a counting model, measured}\newline \textbf{III.3.1~no joint event along the divisors, and the power of the test}\newline \textbf{III.4.1~no deviation from equidistribution, memory or coupling}\newline \emph{also:} III.1.5~common factors of pairs \\
 & heuristic & III.5.3~a conjecture on the most frequent gap \\
 & open & \emph{the most frequent gap} (Problem~\problemvon{III:q:luecken}); \emph{powerful values of $x^4 + 1$} (Problem~\problemvon{III:q:phiacht}) \\
\hline
\end{longtable}
\endgroup
